\section{Fast Apply：把模型输出变成可验证编辑}

模型生成 patch 后，系统仍要定位文件位置、处理用户同时修改、应用 diff、保持格式并触发 diagnostics。Cursor 后来的 Fast Apply 公开资料说明 specialized model 可以承担高吞吐 edit application；本节将其作为 durable design pattern，而不是把后续性能倒写进课堂。

\subsection{Plan、generate、apply、validate}

上一章解决“从哪里获得模型输出”，本章解决“怎样安全进入 editor state”。读图时重点看 plan 与 apply 的职责分离，以及 validation/undo 为什么是产品基础设施。若 apply 静默修改错误位置，即使文本质量高也会破坏信任。

\lecturefigure{13-fast-apply.png}{Planning model 与 specialized apply model 分担理解和高速编辑任务。}{本地字幕 04:20--05:20；Cursor Fast Apply 官方资料。}

\begin{knowledgebox}{读图：Patch application 的四个失败面}
定位错误、context drift、语法/类型失败、用户并发编辑。系统应使用 file/version precondition、structured edit、parser/diagnostic、visible diff 和 atomic undo。对大改动，还需要分批 apply 与中间 validation，避免 editor 被单次巨型 patch 锁死。
\end{knowledgebox}

\begin{importantbox}{交互速度不是单纯 tokens/s}
用户感知时间包含 first useful preview、edit stabilization 与 validation complete。一个更慢但可渐进展示、冲突少、能快速 undo 的 apply path，可能比纯 token throughput 更好。指标应包括 accepted edit、rework 和 rollback，而不仅是生成速度。
\end{importantbox}

\subsection{本章小结}

Fast Apply 把 AI coding 从文本生成问题变成 state transition 问题。模型负责提出修改，系统必须负责 precondition、application、validation、visibility 与 recovery。

